\section{Conservativeness by construction}
\label{sec:guarantee}

The measurements in \cref{sec:results} report no missed collision and no late
time of impact over eleven million query-mode checks. That is evidence and not a
proof, and on its own it leaves open the possibility of an untested query that
fails. This section gives the structural argument, which is short. It then says
exactly where an implementation can break that argument, which in our experience
is the part that gets used.

\begin{theorem}[Conservativeness]
\label{thm:conservative}
Suppose \cref{alg:np} is instantiated so that
\begin{enumerate}
\item[(i)] the rejection test discards a box only when the computed hull of the
  eight corner values, widened by \emph{at least} the certified bound $\epsilon$
  of \cref{eq:err}, excludes the origin on some component; a wider pad is sound,
  and the device kernels use $\max(\tau,\epsilon)$;
\item[(ii)] acceptance reports the box's lower bound in $t$; and
\item[(iii)] every terminating condition other than rejection, whether a depth
  cap, an interval the arithmetic can no longer bisect, or a full work list,
  accepts the box; and
\item[(iv)] the children produced by \op{Split} tile their parent, so every point
  of a parent lies in some child.
\end{enumerate}
Then the reported $\toi$ satisfies $\toi \le \toitrue$, and $\toi < \infty$
whenever a contact occurs.
\end{theorem}

\begin{proof}
Let $\dom$ be any box containing a root of $\Fdiff$. By \cref{eq:hull} the exact
hull of the corner values contains $\Fdiff(\dom) \ni 0$ on every component, and by
\cref{eq:err} the computed hull differs from the exact one by at most $\epsilon$
componentwise. Widening the computed hull by $\epsilon$ therefore contains the
origin, and (i) does not discard $\dom$. So no box containing a root is ever
rejected. By (iv) a box holding the earliest root has a child that also holds it,
so by induction the search retains such a box at every level it reaches.

Consider the terminal box $\dom^{\dagger}$ on that chain, the one at which the
search stops refining. By (iii) it is not dropped, so it is accepted, and by
(ii) the search reports $\dom^{\dagger}.t_{\ell}$. Since $\dom^{\dagger}$
contains the earliest root, $\dom^{\dagger}.t_{\ell} \le \toitrue$. Pruning
against the running minimum only skips boxes whose $t$ lower bound already
exceeds a value the search has attained, which by the same argument is itself
at or below $\toitrue$; so pruning cannot raise the reported value. Hence $\toi
\le \toitrue < \infty$. \end{proof}

Three consequences follow, and each is routinely stated the wrong way round.

\paragraph{Accepting is always safe, however loose the test.} Nothing in
\cref{thm:conservative} constrains \textsc{Accept}. A looser test stops refining
sooner, so it accepts an \emph{ancestor} of the box a tighter test would have
accepted. An ancestor's lower bound in $t$ is at or before its descendant's, so
the looser test reports at or earlier than the tighter one. That costs accuracy
and produces false positives. It cannot produce a late answer or a miss.
\textsc{Relaxed} and \textsc{Tight} therefore differ in accuracy and speed and
not in safety, and the two can be compared on accuracy with no safety caveat
attached.

\paragraph{A depth or stack cap is an accuracy limit, not a safety one.} By
(iii) exhaustion accepts at the box's $t$ lower bound. An undersized cap
therefore reports a time of impact earlier than the true one, and may report a
contact where none exists. It \emph{cannot} miss a collision.

The distinction matters because the effect is large enough to be mistaken for a
correctness failure. A work list sized from the average is far too small for the
tail of \cref{tab:difficulty}. The host explores about $1.2$ boxes per query when
a shared bound is pruning, and about $11$ when it is not, which suggests a
per-thread stack of a few dozen entries. On a vertex-quad query whose exact
answer is $2/3$, that choice costs four orders of magnitude of accuracy:
\begin{center}\ttfamily
\begin{tabular}{ll}
stack cap 32  & 0.6655369599 \\
stack cap 128 & 0.6666666238 \quad\rm(the host's answer)\\
stack cap 512 & 0.6666666238 \\
\end{tabular}
\end{center}
At $32$ entries the answer falls $1.13\times10^{-3}$ before the true value, where
the triangle path reaches $3.0\times10^{-7}$. Every one of the three values is
early, so the invariant holds throughout, and only the step size a solver can
take is lost. Diagnosing this as a conservativeness failure would send an
implementer to the acceptance test, which is not where the problem lives.

The remedy is to derive the capacity. Each pop pushes at
most $k$ children, so a depth-$D$ search needs $1 + kD$ entries. Sizing the work
list from the branching factor and the depth limit makes the tightness of the
answer a property of the search depth. Headroom costs almost nothing: a
$257$-entry stack is a $14.7\,\mathrm{KB}$ frame of which only the first few
entries are ever touched, and register count and spill are unchanged between
capacities $4$ and $513$.

\paragraph{The only way to lose a root is an unsound rejection.} Condition (i)
is the whole of the exposure, and it fails in exactly one way: padding the
origin-containment test by something that merely resembles the certified bound.
Two such substitutions occur in practice, and both are silent.

The first is to pad by the caller's distance tolerance $\tau$ in place of
$\epsilon$. In double precision $\epsilon$ is at most $6.7\times10^{-15}$, while
a typical tolerance is $3\times10^{-8}$. The substitution therefore changes no
answer on any realistic input, and no benchmark detects it. It leaves the search
unsound for any caller who asks for a tolerance below the error bound, and that
is the caller who most needs the guarantee to hold. Padding by
$\max(\tau,\epsilon)$ removes the dependence on what the caller passes. The
second substitution is to pad by the machine epsilon itself, which for a
bilinear query is about $30\times$ too small.

Condition (iii) admits a third failure, and it is the one that most resembles a
paradox. A depth cap that \emph{drops} an exhausted box misses a vertex falling straight
through a stationary triangle in \emph{double} precision, on a
query the same search resolves in single precision. In double the box goes on
subdividing until the cap fires, and the cap then discards it. Exhaustion has to
accept.

Taken together, the three consequences make \cref{thm:conservative} usable as a
checklist. It reduces ``is this implementation conservative'' to three local
questions, and they are the same three for whatever search is being examined.
Does the rejection test pad by the certified bound, or by something that merely
resembles it? On acceptance, is the reported value the box's $t$ lower bound?
On depth, resolution or queue exhaustion, is the box accepted and not dropped?
If all three hold, the kernel is conservative however loose or tight its
acceptance test is, and any measured violation points at one of the three and
not at the algorithm.
